\documentclass[aspectratio=169]{beamer}
\usetheme{Madrid}
\usecolortheme{default}

\usepackage{amsmath,amssymb,amsfonts,bm,mathtools}
\usepackage{physics}
\usepackage{dsfont}

\title{Autoencoders and Representation Learning}
\author{Morten Hjorth-Jensen}
\date{Spring 2026}

\begin{document}

\frame{\titlepage}

\begin{frame}{Course structure}
This mini-course is organized as four one-hour lectures:
\begin{enumerate}
\item Foundations of representation learning and autoencoders
\item Linear autoencoders and the full connection to PCA
\item Nonlinear and regularized autoencoders
\item Probabilistic autoencoders, PPCA, and VAEs
\end{enumerate}

The emphasis is on:
\begin{itemize}
\item linear algebra and optimization,
\item geometry of latent representations,
\item reconstruction error minimization,
\item the relation between autoencoders, PCA, and probabilistic PCA.
\end{itemize}
\end{frame}

\begin{frame}{Notation}
We denote:
\begin{itemize}
\item \(x \in \mathbb{R}^d\): input vector,
\item \(z \in \mathbb{R}^k\): latent code, with \(k<d\) in the undercomplete case,
\item \(X=[x_1,\dots,x_N]\in \mathbb{R}^{d\times N}\): data matrix,
\item \(\hat x\): reconstruction,
\item \(\|A\|_F^2 = \operatorname{Tr}(A^T A)\): Frobenius norm.
\end{itemize}
We will assume throughout that the data are centered when discussing PCA:
\[
\frac{1}{N}\sum_{n=1}^N x_n = 0.
\]
\end{frame}
\section{Lecture 1: Foundations of Autoencoders}

\begin{frame}{What is representation learning?}

Given observations \(x\in\mathbb{R}^d\), representation learning seeks a transformation
\[
x \mapsto z = f(x)
\]
such that \(z\) captures the structure relevant for a downstream task or for reconstruction.

Typical goals:
\begin{itemize}
\item reduce dimension,
\item denoise,
\item identify latent factors,
\item learn a coordinate system adapted to the data.
\end{itemize}

Autoencoders realize this idea through an encoder--decoder architecture trained end-to-end.

\end{frame}

\begin{frame}{Dimensionality reduction as approximation}

A reduced representation of dimension \(k<d\) is useful only if the original data can be approximately recovered:
\[
x \approx \hat x = g(f(x)).
\]
Thus dimensionality reduction is naturally posed as an approximation problem:
\[
\min_{f,g} \sum_{n=1}^N \ell\bigl(x_n,g(f(x_n))\bigr).
\]
The simplest choice is squared error,
\[
\ell(x,\hat x)=\|x-\hat x\|_2^2.
\]
This makes autoencoders closely related to least-squares approximation and low-rank matrix approximation.

\end{frame}

\begin{frame}{Encoder--decoder architecture}

An autoencoder consists of:
\[
z = f_\theta(x), \qquad \hat x = g_\phi(z).
\]
Hence the full model is
\[
\hat x = g_\phi(f_\theta(x)).
\]
Interpretation:
\begin{itemize}
\item \(f_\theta\) compresses data into latent coordinates,
\item \(g_\phi\) reconstructs the original data from those coordinates,
\item the pair \((f_\theta,g_\phi)\) is trained to minimize reconstruction loss.
\end{itemize}

\end{frame}

\begin{frame}{Undercomplete, overcomplete, and bottlenecks}

The latent dimension \(k\) plays a fundamental role.

\textbf{Undercomplete autoencoder:}
\[
k<d.
\]
The network is forced to compress.

\textbf{Overcomplete autoencoder:}
\[
k\ge d.
\]
Without additional regularization, the network may learn the identity map.

Therefore a bottleneck alone is a structural regularizer, but not the only possible one.

\end{frame}

\begin{frame}{Empirical risk for autoencoders}

Given data matrix
\[
X = [x_1,\dots,x_N] \in \mathbb{R}^{d\times N},
\]
the empirical reconstruction objective is
\[
\mathcal{L}(\theta,\phi)
=
\frac{1}{N}\sum_{n=1}^N \|x_n-g_\phi(f_\theta(x_n))\|_2^2.
\]
Equivalently, if \(\hat X=[\hat x_1,\dots,\hat x_N]\),
\[
\mathcal{L}(\theta,\phi)=\frac{1}{N}\|X-\hat X\|_F^2.
\]
Thus the geometry of the Frobenius norm is central from the very beginning.

\end{frame}

\begin{frame}{The Frobenius norm viewpoint}

Recall
\[
\|A\|_F^2 = \sum_{i,j} a_{ij}^2 = \operatorname{Tr}(A^T A).
\]
Hence the reconstruction objective can be written as
\[
\|X-\hat X\|_F^2
=
\operatorname{Tr}\bigl((X-\hat X)^T(X-\hat X)\bigr).
\]
This form is especially useful for deriving optimal linear autoencoders and for proving the connection to PCA.

\end{frame}

\begin{frame}{Latent variables as coordinates}

If the data lie near a \(k\)-dimensional manifold \(\mathcal{M}\subset\mathbb{R}^d\), the encoder can be interpreted as producing local coordinates:
\[
z=f_\theta(x)\in \mathbb{R}^k.
\]
The decoder then acts as a chart map:
\[
g_\phi: \mathbb{R}^k \to \mathbb{R}^d.
\]
Thus autoencoders can be viewed as learning a nonlinear parameterization of the data manifold.

\end{frame}

\begin{frame}{A local linearization viewpoint}

Suppose \(f\) and \(g\) are differentiable. Around a point \(x_0\),
\[
f(x) \approx f(x_0) + J_f(x_0)(x-x_0),
\]
\[
g(z) \approx g(z_0) + J_g(z_0)(z-z_0).
\]
Hence, locally,
\[
g(f(x)) \approx g(f(x_0)) + J_g(z_0)J_f(x_0)(x-x_0).
\]
Thus even nonlinear autoencoders behave like linear dimension reducers in a local neighborhood.

\end{frame}

\begin{frame}{Connection to classical methods}

Autoencoders are related to:
\begin{itemize}
\item PCA: linear projection minimizing squared reconstruction error,
\item factor analysis / probabilistic PCA: latent Gaussian linear models,
\item manifold learning: nonlinear low-dimensional structure,
\item dictionary learning: sparse representations over learned atoms.
\end{itemize}
The key difference is that autoencoders learn both representation and reconstruction maps directly from data.

\end{frame}

\begin{frame}{When does reconstruction imply representation?}

A low reconstruction error does not automatically imply a ``good'' representation in every sense.

Questions to keep in mind:
\begin{itemize}
\item Does the latent space separate classes?
\item Is the representation stable under perturbations?
\item Is the map approximately invertible on the data manifold?
\item Does the latent code capture semantic or physical variables?
\end{itemize}
This motivates adding regularization beyond plain reconstruction.

\end{frame}

\begin{frame}{A general regularized objective}

A broad class of autoencoder objectives has the form
\[
\mathcal{L}(\theta,\phi)
=
\frac{1}{N}\sum_{n=1}^N \|x_n-g_\phi(f_\theta(x_n))\|_2^2
+
\lambda\,\Omega(\theta,\phi;X),
\]
where \(\Omega\) may penalize:
\begin{itemize}
\item latent activity (sparse AE),
\item sensitivity of the encoder (contractive AE),
\item deviation from a prior (VAE),
\item complexity of the weights.
\end{itemize}

\end{frame}

\begin{frame}{Why linear autoencoders matter}

Although deep nonlinear autoencoders are more expressive, linear autoencoders are foundational because:
\begin{itemize}
\item they admit a full analytic treatment,
\item they reveal the relation between reconstruction and projection,
\item they recover PCA exactly under mild assumptions.
\end{itemize}
The next lecture develops this connection rigorously.

\end{frame}

\begin{frame}{Optimization problem in abstract form}

Let \(\mathcal{F}\) be a class of encoders and \(\mathcal{G}\) a class of decoders. Then we solve
\[
\min_{f\in\mathcal{F},\, g\in\mathcal{G}}
\sum_{n=1}^N \|x_n-g(f(x_n))\|_2^2.
\]
If \(\mathcal{F}\) and \(\mathcal{G}\) are linear, the problem becomes a constrained quadratic optimization problem.
If they are nonlinear neural networks, the problem becomes nonconvex but highly expressive.

\end{frame}

\begin{frame}{Autoencoders as identity approximation with constraints}

An autoencoder attempts to approximate the identity map
\[
x \mapsto x,
\]
but only through the restricted factorization
\[
x \mapsto z=f(x)\mapsto \hat x=g(z).
\]
Hence the entire method can be understood as:
\begin{quote}
Find the best constrained approximation to the identity on the data cloud.
\end{quote}
When the constraint is low rank and linear, the optimal solution is PCA.

\end{frame}

\begin{frame}{Reconstruction error decomposition (preview)}

A recurring theme will be:
\[
\|X-\hat X\|_F^2 = \text{total variance} - \text{explained variance}.
\]
For PCA and linear autoencoders, minimizing reconstruction error is equivalent to maximizing variance captured by a \(k\)-dimensional subspace.
This identity will become explicit in Lecture 2.

\end{frame}

\begin{frame}{Lecture 1 summary}

Main points:
\begin{itemize}
\item Autoencoders learn compressed latent representations through reconstruction.
\item The reconstruction objective is naturally expressed in Frobenius norm.
\item Linear autoencoders form the bridge to PCA.
\item Nonlinear and regularized autoencoders extend the idea beyond linear subspaces.
\end{itemize}

\end{frame}

\section{Lecture 2: Linear Autoencoders and PCA}

\begin{frame}{Linear autoencoder architecture}

A linear autoencoder uses
\[
z = W_e x,\qquad \hat x = W_d z = W_d W_e x,
\]
with
\[
W_e\in\mathbb{R}^{k\times d},\qquad W_d\in\mathbb{R}^{d\times k}.
\]
Thus the reconstruction map is linear:
\[
\hat x = A x,\qquad A=W_dW_e.
\]
Since \(A\) factors through \(\mathbb{R}^k\),
\[
\rank(A)\le k.
\]

\end{frame}

\begin{frame}{Matrix formulation of the training problem}

For centered data \(X=[x_1,\dots,x_N]\in\mathbb{R}^{d\times N}\), the objective is
\[
\min_{W_d,W_e}\|X-W_dW_eX\|_F^2.
\]
Defining \(A=W_dW_e\), we rewrite this as
\[
\min_{\rank(A)\le k}\|X-AX\|_F^2.
\]
Therefore linear autoencoder training is a rank-constrained approximation problem.

\end{frame}

\begin{frame}{Sample covariance matrix}

Assume
\[
\frac{1}{N}\sum_{n=1}^N x_n = 0.
\]
Then the empirical covariance matrix is
\[
S = \frac{1}{N}XX^T.
\]
It is symmetric and positive semidefinite, so
\[
S=U\Lambda U^T,
\]
where
\[
\Lambda = \diag(\lambda_1,\dots,\lambda_d),
\qquad
\lambda_1\ge \lambda_2\ge \cdots \ge \lambda_d\ge 0.
\]

\end{frame}

\begin{frame}{PCA as variance maximization}

The principal \(k\)-dimensional subspace is spanned by the top \(k\) eigenvectors:
\[
U_k = [u_1,\dots,u_k].
\]
PCA solves
\[
\max_{B^TB=I_k}\Tr(B^T S B).
\]
Interpretation:
\begin{itemize}
\item columns of \(B\) form an orthonormal basis for a \(k\)-dimensional subspace,
\item \(\Tr(B^TSB)\) is the variance captured by that subspace.
\end{itemize}

\end{frame}

\begin{frame}{PCA as projection minimization}

Equivalent formulation:
\[
\min_{P^2=P,\; P^T=P,\; \rank(P)=k}\|X-PX\|_F^2.
\]
The optimal projector is
\[
P^\star = U_k U_k^T.
\]
Thus the PCA reconstruction is
\[
\hat x = U_k U_k^T x.
\]

\end{frame}

\begin{frame}{Expanding the linear AE loss}

Let \(A\in\mathbb{R}^{d\times d}\) with \(\rank(A)\le k\). Then
\[
\|X-AX\|_F^2
=
\Tr\bigl((X-AX)^T(X-AX)\bigr).
\]
Expanding:
\[
\|X-AX\|_F^2
=
\Tr(X^TX)-2\Tr(X^TAX)+\Tr(X^TA^TAX).
\]
Using cyclicity of the trace and \(XX^T = NS\),
\[
\|X-AX\|_F^2
=
N\Bigl[\Tr(S)-2\Tr(AS)+\Tr(A^TAS)\Bigr].
\]

\end{frame}

\begin{frame}{Reduction to orthogonal projection}

Let \(\mathcal{M}=\operatorname{Im}(A)\), with \(\dim \mathcal{M}\le k\).
For each \(x_n\), the reconstruction \(Ax_n\in \mathcal{M}\).
Let \(P_{\mathcal{M}}\) be the orthogonal projector onto \(\mathcal{M}\). Then
\[
\|x_n-P_{\mathcal{M}}x_n\|_2^2 \le \|x_n-Ax_n\|_2^2
\]
for every \(n\), because orthogonal projection is the closest point in the subspace.

Therefore
\[
\|X-P_{\mathcal{M}}X\|_F^2 \le \|X-AX\|_F^2.
\]
So the global optimum may be taken to be an orthogonal projector of rank at most \(k\).

\end{frame}

\begin{frame}{Projector parameterization}

Any rank-\(k\) orthogonal projector can be written as
\[
P = BB^T,
\]
where
\[
B\in\mathbb{R}^{d\times k},\qquad B^TB=I_k.
\]
Then
\[
\|X-PX\|_F^2
=
N\Bigl[\Tr(S)-\Tr(B^TSB)\Bigr].
\]
Thus minimizing the reconstruction error is equivalent to maximizing
\[
\Tr(B^TSB)
\quad \text{subject to} \quad B^TB=I_k.
\]

\end{frame}

\begin{frame}{Variational characterization in eigenbasis}

Let
\[
S = U\Lambda U^T,\qquad C = U^T B.
\]
Since \(U\) is orthogonal and \(B^TB=I_k\),
\[
C^T C = I_k.
\]
Now
\[
\Tr(B^TSB) = \Tr(C^T \Lambda C).
\]
Write \(c_{ij}\) for the entries of \(C\). Then
\[
\Tr(C^T\Lambda C)=\sum_{i=1}^d \lambda_i \alpha_i,
\qquad
\alpha_i = \sum_{j=1}^k c_{ij}^2.
\]
Also,
\[
0\le \alpha_i \le 1,\qquad \sum_{i=1}^d \alpha_i = k.
\]

\end{frame}

\begin{frame}{Maximization argument}

Because
\[
\lambda_1\ge \lambda_2\ge \cdots \ge \lambda_d,
\]
the quantity
\[
\sum_{i=1}^d \lambda_i \alpha_i
\]
is maximized by placing all available weight on the largest \(k\) eigenvalues:
\[
\alpha_1=\cdots=\alpha_k=1,\qquad \alpha_{k+1}=\cdots=\alpha_d=0.
\]
Hence
\[
\max_{B^TB=I_k}\Tr(B^TSB)=\sum_{i=1}^k \lambda_i,
\]
and the maximizing subspace is
\[
\operatorname{span}(B)=\operatorname{span}(u_1,\dots,u_k).
\]

\end{frame}

\begin{frame}{Optimal reconstruction map}

Therefore the optimal rank-\(k\) projector is
\[
P^\star = U_k U_k^T.
\]
Since the optimal linear autoencoder may be chosen to be a projector,
\[
A^\star = W_d^\star W_e^\star = U_k U_k^T.
\]
So linear autoencoders recover the PCA subspace.

\end{frame}

\begin{frame}{Choice of encoder and decoder}

A canonical choice is
\[
W_e^\star = U_k^T,\qquad W_d^\star = U_k.
\]
Then
\[
W_d^\star W_e^\star = U_k U_k^T.
\]
The latent code is
\[
z = U_k^T x,
\]
which is exactly the vector of principal coordinates, and the reconstruction is
\[
\hat x = U_k z = U_k U_k^T x.
\]

\end{frame}

\begin{frame}{Non-uniqueness of the factors}

The product \(A^\star=U_kU_k^T\) is determined by the principal subspace, but \(W_e^\star\) and \(W_d^\star\) are not unique.
For any invertible \(M\in\mathbb{R}^{k\times k}\),
\[
W_d = U_k M,\qquad W_e = M^{-1} U_k^T
\]
also satisfy
\[
W_d W_e = U_kU_k^T.
\]
Hence the latent coordinates are unique only up to an invertible linear change of basis.

\end{frame}

\begin{frame}{Reconstruction error and discarded variance}

The optimal reconstruction error is
\[
\|X-U_kU_k^T X\|_F^2
=
N\sum_{i=k+1}^d \lambda_i.
\]
So PCA and the optimal linear autoencoder discard exactly the variance associated with the trailing eigenvalues.
This is one of the most useful interpretations of PCA:
\[
\text{reconstruction error} = \text{unexplained variance}.
\]

\end{frame}

\begin{frame}{Connection to SVD}

Let
\[
X = U\Sigma V^T
\]
be the singular value decomposition.
Then
\[
S = \frac{1}{N}XX^T = \frac{1}{N}U\Sigma^2 U^T.
\]
Hence the PCA directions are the left singular vectors of \(X\).
The optimal linear autoencoder reconstructs with
\[
A^\star = U_k U_k^T,
\]
where \(U_k\) contains the first \(k\) left singular vectors.

\end{frame}

\begin{frame}{Relation to Eckart--Young--Mirsky}

The Eckart--Young--Mirsky theorem states that the best rank-\(k\) approximation to \(X\) in Frobenius norm is
\[
X_k = U_k \Sigma_k V_k^T.
\]
The projector form is related by
\[
U_k U_k^T X = U_k U_k^T (U\Sigma V^T)=U_k\Sigma_k V_k^T = X_k.
\]
Thus the linear autoencoder recovers the same approximation as truncated SVD.

\end{frame}

\begin{frame}{Bias terms and centering}

If biases are included,
\[
z = W_e x + b_e,\qquad \hat x = W_d z + b_d,
\]
then the relation to PCA is cleanest after centering the data.
Otherwise the affine part tends to learn the empirical mean.
In practice:
\begin{itemize}
\item center first for the pure PCA theorem,
\item or interpret the bias as handling the data mean.
\end{itemize}

\end{frame}

\begin{frame}{Orthonormality and tied weights}

Sometimes one imposes tied weights:
\[
W_d = W_e^T.
\]
Then the model becomes
\[
\hat x = W_e^T W_e x.
\]
If the rows of \(W_e\) are orthonormal, then \(W_e^TW_e\) is the orthogonal projector onto the row space of \(W_e\).
This makes the PCA connection even more direct, although tied weights are not necessary for the theorem.

\end{frame}

\begin{frame}{Geometric meaning of the theorem}

Linear autoencoders do not merely learn a low-dimensional code:
they learn the \emph{best} \(k\)-dimensional linear subspace in the least-squares sense.

So the theorem says:
\begin{quote}
Among all \(k\)-dimensional linear reconstructions, the optimal one is orthogonal projection onto the principal subspace.
\end{quote}
That is exactly the geometric content of PCA.

\end{frame}

\begin{frame}{Lecture 2 summary}

Main results:
\begin{itemize}
\item The linear autoencoder objective is a rank-constrained least-squares problem.
\item Its global minimizer is an orthogonal projector.
\item The optimal projector is \(U_kU_k^T\), the PCA projector.
\item Therefore linear autoencoders recover PCA exactly.
\end{itemize}

\end{frame}

\section{Lecture 3: Nonlinear and Regularized Autoencoders}

\begin{frame}{Why go beyond PCA?}

PCA captures only linear subspaces.
But many datasets lie near nonlinear manifolds:
\begin{itemize}
\item images of a rotated object,
\item molecular conformations,
\item solutions of nonlinear PDEs,
\item speech signals.
\end{itemize}
A nonlinear autoencoder replaces linear maps by compositions of affine maps and nonlinear activations.

\end{frame}

\begin{frame}{Deep autoencoder architecture}

A deep encoder may have the form
\[
z = f_\theta(x)=\sigma_L(W_L \sigma_{L-1}(W_{L-1}\cdots \sigma_1(W_1 x+b_1)\cdots+b_{L-1})+b_L),
\]
and the decoder similarly
\[
\hat x = g_\phi(z).
\]
These nested nonlinear maps can approximate highly non-Euclidean structure.

\end{frame}

\begin{frame}{Local linearization of nonlinear autoencoders}

If \(f\) and \(g\) are differentiable, then around \(x_0\),
\[
f(x) \approx f(x_0)+J_f(x_0)(x-x_0),
\]
\[
g(z) \approx g(z_0)+J_g(z_0)(z-z_0).
\]
Hence
\[
g(f(x)) \approx g(f(x_0)) + J_g(z_0)J_f(x_0)(x-x_0).
\]
So the Jacobian product \(J_g J_f\) acts as the local linear reconstruction operator.

\end{frame}

\begin{frame}{Latent manifolds and charts}

If the data lie on a manifold \(\mathcal{M}\subset\mathbb{R}^d\) of intrinsic dimension \(k\), then
\[
f_\theta : \mathcal{M}\to\mathbb{R}^k
\]
can be interpreted as a chart, and
\[
g_\phi : \mathbb{R}^k\to\mathcal{M}
\]
as a reconstruction map.

Autoencoders therefore fit naturally into differential-geometric language:
\begin{itemize}
\item latent space gives coordinates,
\item decoder parameterizes the manifold in ambient space.
\end{itemize}

\end{frame}

\begin{frame}{Denoising autoencoders}

In a denoising autoencoder, one corrupts the input:
\[
\tilde x = x + \varepsilon,
\]
and trains the network to reconstruct \(x\) from \(\tilde x\):
\[
\min_{\theta,\phi}\frac{1}{N}\sum_{n=1}^N \|x_n-g_\phi(f_\theta(\tilde x_n))\|_2^2.
\]
This prevents the model from trivially copying the input and encourages extraction of robust structure.

\end{frame}

\begin{frame}{Statistical meaning of denoising}

For small Gaussian noise, the denoising map approximates a score-like direction.
In fact, denoising autoencoders are closely related to score matching:
the reconstruction vector
\[
r(x)-x
\]
encodes information about the gradient of the log-density of the data distribution.
Thus denoising autoencoders anticipate ideas that later appear in diffusion and score-based models.

\end{frame}

\begin{frame}{Sparse autoencoders}

Sparse autoencoders add a penalty encouraging latent activations to be mostly zero:
\[
\mathcal{L} =
\frac{1}{N}\sum_{n=1}^N \|x_n-\hat x_n\|_2^2
+\lambda \sum_{n=1}^N \|z_n\|_1.
\]
Interpretation:
\begin{itemize}
\item each input activates only a small subset of latent units,
\item latent units can become feature detectors,
\item the model resembles dictionary learning and sparse coding.
\end{itemize}

\end{frame}

\begin{frame}{Contractive autoencoders}

Contractive autoencoders penalize the sensitivity of the encoder:
\[
\mathcal{L} =
\frac{1}{N}\sum_{n=1}^N \|x_n-\hat x_n\|_2^2
+
\lambda \sum_{n=1}^N \|J_f(x_n)\|_F^2.
\]
This encourages the representation to be stable under small perturbations orthogonal to the data manifold.
Hence the encoder becomes approximately invariant to directions that do not matter.

\end{frame}

\begin{frame}{Jacobian penalty and geometry}

The Jacobian \(J_f(x)\) measures how latent coordinates change with the input.
If
\[
\|J_f(x)\|_F^2
\]
is small in directions orthogonal to the manifold, then the encoder contracts those directions.
So contractive autoencoders can be seen as learning tangent directions of the data manifold while suppressing normal directions.

\end{frame}

\begin{frame}{Overcomplete autoencoders and regularization}

An overcomplete autoencoder with \(k\ge d\) can easily learn
\[
g(f(x))=x.
\]
To prevent this, one must impose additional structure:
\begin{itemize}
\item sparsity,
\item denoising corruption,
\item contractive penalty,
\item weight decay or architectural constraints.
\end{itemize}
Thus regularization can replace the role of a bottleneck.

\end{frame}

\begin{frame}{Kernel PCA vs nonlinear autoencoders}

Kernel PCA uses a fixed nonlinear feature map \(\Phi(x)\) and performs PCA in feature space.
A nonlinear autoencoder instead learns both:
\begin{itemize}
\item the representation map,
\item and the inverse reconstruction map.
\end{itemize}
Kernel PCA is nonparametric and spectral;
autoencoders are parametric and optimization-based.

\end{frame}

\begin{frame}{Relation to manifold learning methods}

Classical manifold learning methods include:
\begin{itemize}
\item Isomap,
\item LLE,
\item Laplacian eigenmaps,
\item diffusion maps.
\end{itemize}
These methods often construct embeddings from graph geometry.
Autoencoders instead learn a parametric map
\[
x \mapsto z
\]
and an explicit decoder
\[
z \mapsto \hat x,
\]
which is often advantageous for out-of-sample extension.

\end{frame}

\begin{frame}{Latent interpolation}

One practical way to test a learned latent geometry is interpolation:
\[
z(\tau)=(1-\tau)z_1+\tau z_2,\qquad 0\le \tau \le 1,
\]
followed by decoding
\[
\hat x(\tau)=g_\phi(z(\tau)).
\]
If the autoencoder has learned a smooth latent manifold, then these interpolations should remain data-like.

\end{frame}

\begin{frame}{Optimization landscape}

Nonlinear autoencoder training is nonconvex because of:
\begin{itemize}
\item multilayer compositions,
\item nonlinear activations,
\item coupled encoder/decoder parameters.
\end{itemize}
Thus unlike linear autoencoders, we generally do not have closed-form global minimizers.
Instead, gradient-based optimization is used.

\end{frame}

\begin{frame}{Generalization issues}

A nonlinear autoencoder can memorize the training set if too expressive.
Hence we must ask:
\begin{itemize}
\item does it generalize to unseen data?
\item is the latent structure stable?
\item does reconstruction preserve physically meaningful variation?
\end{itemize}
These questions motivate probabilistic and information-theoretic viewpoints.

\end{frame}

\begin{frame}{Lecture 3 summary}

Main ideas:
\begin{itemize}
\item Nonlinear autoencoders generalize PCA from linear subspaces to nonlinear manifolds.
\item Denoising, sparse, and contractive variants add structure to the latent representation.
\item Jacobians provide the local differential geometry of the encoder and decoder.
\end{itemize}

\end{frame}

\section{Lecture 4: Probabilistic Autoencoders, PPCA, and VAEs}

\begin{frame}{From deterministic to probabilistic models}

Deterministic autoencoders produce a code
\[
z=f_\theta(x)
\]
and a reconstruction
\[
\hat x=g_\phi(z).
\]
A probabilistic autoencoder instead introduces a latent-variable model
\[
p_\theta(x,z)=p_\theta(x\mid z)p(z),
\]
where \(z\) is a random latent variable.
This turns representation learning into statistical inference.

\end{frame}

\begin{frame}{Probabilistic PCA model}

Probabilistic PCA assumes
\[
z \sim \mathcal{N}(0,I_k),
\]
and
\[
x = W z + \mu + \epsilon,
\qquad
\epsilon \sim \mathcal{N}(0,\sigma^2 I_d).
\]
Hence
\[
p(x\mid z)=\mathcal{N}(Wz+\mu,\sigma^2 I_d).
\]
This is a linear-Gaussian latent variable model.

\end{frame}

\begin{frame}{Marginal distribution in PPCA}

Integrating out \(z\), one gets
\[
x \sim \mathcal{N}(\mu,\, WW^T+\sigma^2 I_d).
\]
Thus PPCA models the covariance by a low-rank term plus isotropic noise:
\[
\Sigma = WW^T + \sigma^2 I_d.
\]
This is the probabilistic analogue of low-dimensional linear structure.

\end{frame}

\begin{frame}{PPCA and classical PCA}

As \(\sigma^2 \to 0\), the PPCA model approaches a deterministic projection model.
Moreover, maximum-likelihood estimation in PPCA yields a loading matrix whose column space is the PCA subspace.
So PPCA recovers the same principal directions as classical PCA, but with a probabilistic interpretation.

\end{frame}

\begin{frame}{Posterior distribution in PPCA}

Because the model is linear-Gaussian, the posterior is analytic:
\[
p(z\mid x)=\mathcal{N}(M^{-1}W^T(x-\mu),\;\sigma^2 M^{-1}),
\qquad
M=W^TW+\sigma^2 I_k.
\]
Thus PPCA has exact inference, unlike general nonlinear latent-variable models.

\end{frame}

\begin{frame}{EM perspective for PPCA}

PPCA can be trained by EM:
\begin{itemize}
\item E-step: compute posterior moments of \(z\) given \(x\),
\item M-step: update \(W\) and \(\sigma^2\) to maximize expected complete-data log-likelihood.
\end{itemize}
This makes PPCA an important bridge between PCA and variational autoencoders.

\end{frame}

\begin{frame}{General latent-variable model}

In a general generative model,
\[
p_\theta(x)=\int p_\theta(x\mid z)p(z)\,dz.
\]
This integral is usually intractable for nonlinear \(p_\theta(x\mid z)\).
So direct maximum likelihood is difficult.
We therefore introduce an approximate posterior \(q_\phi(z\mid x)\).

\end{frame}

\begin{frame}{Variational inference and the ELBO}

Starting from
\[
\log p_\theta(x)
=
\log \int q_\phi(z\mid x)\frac{p_\theta(x,z)}{q_\phi(z\mid x)}\,dz,
\]
Jensen's inequality gives
\[
\log p_\theta(x)
\ge
\mathbb{E}_{q_\phi(z\mid x)}
\left[
\log \frac{p_\theta(x,z)}{q_\phi(z\mid x)}
\right].
\]
This lower bound is the evidence lower bound (ELBO).

\end{frame}

\begin{frame}{ELBO decomposition}

Expanding the ELBO,
\[
\mathcal{L}_{\mathrm{ELBO}}(x)
=
\mathbb{E}_{q_\phi(z\mid x)}[\log p_\theta(x\mid z)]
-
\mathrm{KL}\bigl(q_\phi(z\mid x)\|p(z)\bigr).
\]
Interpretation:
\begin{itemize}
\item reconstruction term encourages faithful decoding,
\item KL term regularizes the latent distribution toward the prior.
\end{itemize}

\end{frame}

\begin{frame}{Why the ELBO is useful}

The ELBO satisfies
\[
\log p_\theta(x)
=
\mathcal{L}_{\mathrm{ELBO}}(x)
+
\mathrm{KL}\bigl(q_\phi(z\mid x)\|p_\theta(z\mid x)\bigr).
\]
Since KL divergence is nonnegative,
\[
\mathcal{L}_{\mathrm{ELBO}}(x)\le \log p_\theta(x).
\]
Maximizing the ELBO therefore simultaneously:
\begin{itemize}
\item increases the data likelihood,
\item makes the approximate posterior \(q_\phi(z\mid x)\) close to the true posterior.
\end{itemize}

\end{frame}

\begin{frame}{Variational autoencoder architecture}

A VAE parameterizes:
\begin{itemize}
\item encoder (approximate posterior): \(q_\phi(z\mid x)\),
\item decoder (likelihood model): \(p_\theta(x\mid z)\),
\item prior: usually \(p(z)=\mathcal{N}(0,I)\).
\end{itemize}
Typically,
\[
q_\phi(z\mid x)=\mathcal{N}(\mu_\phi(x),\diag(\sigma_\phi(x)^2)).
\]
Thus the encoder outputs the parameters of a distribution, not a single deterministic code.

\end{frame}

\begin{frame}{Reparameterization trick}

To differentiate through stochastic sampling, write
\[
z = \mu_\phi(x) + \sigma_\phi(x)\odot \epsilon,
\qquad
\epsilon\sim\mathcal{N}(0,I).
\]
Then the randomness is isolated in \(\epsilon\), while \(\mu_\phi\) and \(\sigma_\phi\) remain differentiable parameters.
This makes stochastic gradient training possible.

\end{frame}

\begin{frame}{VAE objective as autoencoder + regularizer}

The VAE objective can be read as
\[
\text{reconstruction quality}
-
\text{latent complexity penalty}.
\]
Hence the VAE is an autoencoder with a probabilistic latent bottleneck.
Unlike deterministic autoencoders, the latent space is explicitly shaped by the prior.

\end{frame}

\begin{frame}{Information bottleneck viewpoint}

The KL term
\[
\mathrm{KL}(q_\phi(z\mid x)\|p(z))
\]
limits how much information about \(x\) is stored in \(z\).
Thus VAEs are closely related to the information bottleneck principle:
\begin{itemize}
\item keep information useful for reconstruction,
\item discard irrelevant detail,
\item regularize the representation toward a simple latent distribution.
\end{itemize}

\end{frame}

\begin{frame}{Rate--distortion interpretation}

The ELBO also resembles a rate--distortion tradeoff:
\begin{itemize}
\item distortion \(\leftrightarrow\) reconstruction error,
\item rate \(\leftrightarrow\) KL cost of encoding information in \(z\).
\end{itemize}
In \(\beta\)-VAEs, one explicitly tunes this tradeoff:
\[
\mathcal{L}_{\beta\text{-VAE}}
=
\mathbb{E}_{q_\phi(z\mid x)}[\log p_\theta(x\mid z)]
-
\beta\, \mathrm{KL}(q_\phi(z\mid x)\|p(z)).
\]

\end{frame}

\begin{frame}{Linear VAE and PPCA}

If the decoder is linear and Gaussian and the encoder family is chosen appropriately, the VAE reduces to a variational treatment of PPCA.
Thus there is a hierarchy:
\[
\text{PCA}
\longrightarrow
\text{Linear AE}
\longrightarrow
\text{PPCA}
\longrightarrow
\text{VAE}.
\]
Each step adds modeling flexibility and probabilistic structure.

\end{frame}

\begin{frame}{Deterministic AE vs VAE}

Deterministic autoencoder:
\[
z=f_\theta(x),\qquad \hat x = g_\phi(z).
\]

Variational autoencoder:
\[
z\sim q_\phi(z\mid x),\qquad x\sim p_\theta(x\mid z).
\]

Difference:
\begin{itemize}
\item AE optimizes reconstruction directly,
\item VAE optimizes a lower bound on likelihood and learns a structured latent distribution.
\end{itemize}

\end{frame}

\begin{frame}{What PCA misses that VAEs capture}

PCA gives:
\begin{itemize}
\item linear coordinates,
\item variance-maximizing subspace,
\item deterministic projection.
\end{itemize}

VAEs additionally provide:
\begin{itemize}
\item nonlinear latent geometry,
\item generative sampling,
\item uncertainty in latent representations,
\item explicit prior structure.
\end{itemize}

\end{frame}

\begin{frame}{Global summary of the mini-course}

We have developed a progression:
\begin{enumerate}
\item Autoencoders as constrained identity approximators.
\item Linear autoencoders as exact recoverers of PCA.
\item Nonlinear autoencoders as manifold learners.
\item Probabilistic autoencoders and VAEs as latent-variable generative models.
\end{enumerate}

This unifies linear algebra, geometry, optimization, and probabilistic modeling.

\end{frame}

\begin{frame}{Possible follow-up topics}

Natural continuations include:
\begin{itemize}
\item denoising score matching and its relation to denoising autoencoders,
\item normalizing flows and invertible autoencoding ideas,
\item diffusion models as denoising-based generative models,
\item information-theoretic analysis of latent variable learning.
\end{itemize}

\end{frame}

\begin{frame}{Representation learning as compression}

Suppose \(x\) is generated from latent factors \(s\in\mathbb{R}^k\) with \(k\ll d\).
Then a good representation \(z\) should satisfy:
\begin{itemize}
\item \(\dim z\) is small,
\item \(z\) is sufficient for reconstruction,
\item \(z\) is stable under small perturbations.
\end{itemize}
This motivates optimization problems of the form
\[
\min_{f,g}\; \mathbb{E}\|x-g(f(x))\|^2
\]
subject to complexity constraints on \(f\) and \(g\).

\end{frame}

\begin{frame}{Why not learn the identity map?}

If the autoencoder is too expressive, then it may learn
\[
g_\phi(f_\theta(x)) = x
\]
for all training examples without extracting useful structure.

This occurs especially when:
\begin{itemize}
\item the latent space is overcomplete,
\item no regularization is used,
\item the network has enough capacity to memorize the dataset.
\end{itemize}

Hence a useful autoencoder must include a mechanism that prevents trivial copying.

\end{frame}

\begin{frame}{Data manifolds and tangent spaces}

If the data lie on a smooth manifold \(\mathcal{M}\subset \mathbb{R}^d\), then at each point \(x\in\mathcal{M}\) there is a tangent space
\[
T_x\mathcal{M}\subset \mathbb{R}^d.
\]
A good encoder should preserve variation along tangent directions and suppress variation normal to the manifold.
This geometric viewpoint explains why Jacobians play a key role in regularized autoencoders.

\end{frame}

\begin{frame}{Loss functions beyond squared error}

Although the standard choice is
\[
\ell(x,\hat x)=\|x-\hat x\|_2^2,
\]
other reconstruction losses are often appropriate:
\begin{itemize}
\item binary cross-entropy for binary or Bernoulli data,
\item weighted quadratic loss for heteroscedastic noise,
\item perceptual losses in image applications.
\end{itemize}
However, squared error is the natural loss for the link to PCA.

\end{frame}

\begin{frame}{Batch optimization form}

For a dataset matrix \(X\), the encoder produces latent matrix \(Z\), and the decoder produces
\[
\hat X = g_\phi(f_\theta(X)).
\]
The batch objective becomes
\[
\mathcal{L} = \frac{1}{N}\|X-\hat X\|_F^2.
\]
This matrix form emphasizes that an autoencoder learns a transformation of an entire cloud of points, not merely independent reconstructions.

\end{frame}

\begin{frame}{Linear algebra of latent codes}

If the latent code for the \(n\)-th sample is \(z_n\), then collecting them gives
\[
Z=[z_1,\dots,z_N]\in \mathbb{R}^{k\times N}.
\]
The decoder then maps \(Z\) back to \(\hat X\in \mathbb{R}^{d\times N}\).
Thus the autoencoder implicitly factorizes the data through a low-dimensional representation:
\[
X \mapsto Z \mapsto \hat X.
\]
This already hints at matrix factorization and low-rank approximation.

\end{frame}

\begin{frame}{Autoencoders and inverse problems}

In many applications, \(x\) is not observed directly but through a forward model \(y=Ax+\eta\).
An autoencoder can then be used as a structural prior:
\[
x \approx g_\phi(z).
\]
Instead of solving for \(x\in\mathbb{R}^d\), one solves for \(z\in\mathbb{R}^k\), which reduces dimension and imposes regularity.
This is one reason autoencoders are useful beyond pure unsupervised learning.

\end{frame}

\begin{frame}{Lecture 1 conceptual exercise}

Show that if \(g\circ f\) is linear and idempotent,
\[
(g\circ f)^2 = g\circ f,
\]
then the reconstruction map is a projector.
Why is this already suggestive of a connection to PCA?

\end{frame}

\begin{frame}{PCA from a constrained basis}

If \(B=[b_1,\dots,b_k]\in\mathbb{R}^{d\times k}\) with \(B^TB=I_k\), then the projector onto \(\operatorname{span}(B)\) is
\[
P_B = BB^T.
\]
The PCA problem can be written as
\[
\min_{B^TB=I_k}\|X-BB^TX\|_F^2.
\]
Thus PCA is a search over orthonormal \(k\)-frames in \(\mathbb{R}^d\).

\end{frame}

\begin{frame}{Trace expansion for PCA}

Using \(P=BB^T\),
\[
\|X-PX\|_F^2
=
\Tr(X^TX)-2\Tr(X^TPX)+\Tr(X^TP^TPX).
\]
Since \(P^T=P\) and \(P^2=P\),
\[
\|X-PX\|_F^2=\Tr(X^TX)-\Tr(X^TPX).
\]
Because \(XX^T=NS\),
\[
\|X-PX\|_F^2 = N\bigl(\Tr(S)-\Tr(B^TSB)\bigr).
\]

\end{frame}

\begin{frame}{Why orthogonal projection is best in a fixed subspace}

Fix a subspace \(\mathcal{M}\subset \mathbb{R}^d\). Among all \(y\in \mathcal{M}\), the orthogonal projection \(P_{\mathcal{M}}x\) uniquely minimizes
\[
\|x-y\|_2^2.
\]
Proof: write \(x = P_{\mathcal{M}}x + r\), with \(r\perp \mathcal{M}\). Then for any \(y\in\mathcal{M}\),
\[
\|x-y\|_2^2
=
\|P_{\mathcal{M}}x-y\|_2^2 + \|r\|_2^2 \ge \|r\|_2^2.
\]
This is the key step allowing us to reduce the linear autoencoder problem to PCA.

\end{frame}

\begin{frame}{Lagrange multiplier derivation of PCA}

To maximize
\[
\Tr(B^TSB)
\quad \text{subject to} \quad B^TB=I_k,
\]
consider
\[
\mathcal{J}(B,\Lambda)=\Tr(B^TSB)-\Tr\bigl(\Lambda(B^TB-I_k)\bigr).
\]
Differentiating with respect to \(B\) gives
\[
2SB - 2B\Lambda = 0,
\]
so
\[
SB = B\Lambda.
\]
Hence the columns of \(B\) are eigenvectors of \(S\).

\end{frame}

\begin{frame}{Explained variance}

For the principal projector \(P_k=U_kU_k^T\),
\[
\Tr(P_k S)=\sum_{i=1}^k \lambda_i.
\]
This is called the explained variance.
The fraction of variance explained is
\[
\frac{\sum_{i=1}^k \lambda_i}{\sum_{i=1}^d \lambda_i}.
\]
This provides a principled way to choose the bottleneck dimension \(k\) in a linear autoencoder.

\end{frame}

\begin{frame}{Orthogonality of reconstruction error}

For PCA projection,
\[
x = U_kU_k^T x + (I-U_kU_k^T)x.
\]
The residual satisfies
\[
U_k^T (I-U_kU_k^T)x = 0.
\]
Thus the reconstruction error is orthogonal to the retained subspace, which is another hallmark of least-squares optimality.

\end{frame}

\begin{frame}{Tied weights and symmetry}

With tied weights \(W_d=W_e^T\), the linear autoencoder objective becomes
\[
\min_{W_e}\|X-W_e^TW_eX\|_F^2.
\]
If \(W_eW_e^T=I_k\), then \(W_e^TW_e\) is a rank-\(k\) orthogonal projector.
Therefore tied-weight linear autoencoders naturally encourage PCA-like structure.

\end{frame}

\begin{frame}{When the PCA subspace is not unique}

If
\[
\lambda_k = \lambda_{k+1},
\]
then the principal \(k\)-dimensional subspace need not be unique.
In this case, linear autoencoders still recover a principal subspace, but not a unique one.
The non-uniqueness comes from both:
\begin{itemize}
\item degeneracy of the eigenspace,
\item invertible changes of latent basis.
\end{itemize}

\end{frame}

\begin{frame}{Centered versus uncentered data}

PCA theory assumes centered data. If \(X\) is not centered, then minimizing
\[
\|X-W_dW_eX\|_F^2
\]
need not coincide exactly with PCA of the covariance matrix.
In practice, one often centers:
\[
X \leftarrow X - \bar x \mathds{1}^T.
\]
Then the linear autoencoder theorem applies directly.

\end{frame}

\begin{frame}{Lecture 2 proof checkpoint}

The full proof that linear autoencoders recover PCA had four key steps:
\begin{enumerate}
\item rewrite the model using \(A=W_dW_e\),
\item reduce to projection onto a \(k\)-dimensional subspace,
\item solve the projection problem by spectral optimization,
\item identify the optimal projector as \(U_kU_k^T\).
\end{enumerate}

\end{frame}

\begin{frame}{Lecture 2 exercise}

Show directly from the SVD \(X=U\Sigma V^T\) that
\[
U_kU_k^TX = U_k\Sigma_k V_k^T.
\]
Why does this imply that the PCA reconstruction is the best rank-\(k\) approximation to \(X\)?

\end{frame}

\begin{frame}{Activation functions and expressivity}

Common activations in autoencoders include:
\begin{itemize}
\item sigmoid,
\item hyperbolic tangent,
\item ReLU and leaky ReLU,
\item smooth activations such as GELU or softplus.
\end{itemize}
The activation determines both expressive power and local geometry through its derivative.

\end{frame}

\begin{frame}{Universal approximation intuition}

A feed-forward network with nonlinear activations can approximate broad classes of continuous maps.
Hence a nonlinear autoencoder can in principle learn a curved manifold embedding
\[
f:\mathbb{R}^d \to \mathbb{R}^k,\qquad g:\mathbb{R}^k\to \mathbb{R}^d
\]
far beyond the linear subspace model of PCA.

\end{frame}

\begin{frame}{Decoder as manifold parameterization}

The decoder map
\[
g_\phi(z)
\]
can be interpreted as a learned surface embedded in \(\mathbb{R}^d\).
If \(k=2\), for example, then \(g_\phi\) parameterizes a 2D surface inside a high-dimensional ambient space.
Thus nonlinear autoencoders replace linear subspaces by learned immersed manifolds.

\end{frame}

\begin{frame}{Encoder Jacobian and sensitivity}

For encoder \(f:\mathbb{R}^d\to\mathbb{R}^k\), the Jacobian is
\[
J_f(x) = \frac{\partial f(x)}{\partial x}\in \mathbb{R}^{k\times d}.
\]
Its singular values quantify local sensitivity:
\begin{itemize}
\item large singular values: latent code changes rapidly,
\item small singular values: local contraction.
\end{itemize}
Regularizers based on \(J_f\) therefore control local geometry.

\end{frame}

\begin{frame}{Denoising as manifold projection}

If noisy observations satisfy
\[
\tilde x = x + \varepsilon,
\]
with \(x\) on or near a manifold \(\mathcal{M}\), then denoising can be interpreted as approximately mapping \(\tilde x\) back toward \(\mathcal{M}\).
So a denoising autoencoder learns a map resembling a noisy manifold projection.

\end{frame}

\begin{frame}{Sparse coding viewpoint}

Sparse autoencoders are closely related to the optimization
\[
\min_{D,Z} \|X-DZ\|_F^2 + \lambda \|Z\|_1,
\]
where \(D\) is a dictionary and \(Z\) contains sparse codes.
Thus sparse autoencoders connect neural representation learning with classical sparse approximation.

\end{frame}

\begin{frame}{Contractive penalty in coordinates}

For a one-hidden-layer encoder
\[
h = \sigma(Wx+b),
\]
the Jacobian is
\[
J_f(x)=\diag(\sigma'(Wx+b))\,W.
\]
Hence the contractive penalty
\[
\|J_f(x)\|_F^2
\]
depends on both weights and activation derivatives, and suppresses sensitivity especially when activations saturate.

\end{frame}

\begin{frame}{Robustness to perturbations}

A desirable latent representation should satisfy
\[
\|f(x+\delta)-f(x)\| \ll \|\delta\|
\]
for nuisance perturbations \(\delta\), while still preserving relevant variation.
Contractive and denoising objectives both encourage such robustness, but in different ways:
\begin{itemize}
\item denoising acts through corrupted inputs,
\item contractive training acts through local derivative control.
\end{itemize}

\end{frame}

\begin{frame}{Nonlinear autoencoders vs PCA locally}

Although nonlinear autoencoders are globally nonlinear, around any point \(x_0\) they induce a local linear map
\[
J_g(z_0)J_f(x_0).
\]
This local map can be compared to the PCA projector.
Thus one can interpret nonlinear autoencoders as learning a data-dependent family of local projections rather than one global projection.

\end{frame}

\begin{frame}{Lecture 3 exercise}

Suppose a denoising autoencoder is trained with small Gaussian noise.
Why should the reconstruction vector \(r(x)-x\) point approximately toward high-density regions of the data distribution?
Relate your reasoning to denoising and local score estimation.

\end{frame}

\begin{frame}{Complete-data log likelihood in PPCA}

For the PPCA model,
\[
p(x,z)=p(x\mid z)p(z),
\]
the complete-data log likelihood is, up to constants,
\[
-\frac{1}{2\sigma^2}\|x-Wz-\mu\|_2^2 - \frac{1}{2}\|z\|_2^2.
\]
This quadratic structure is what makes EM and exact posterior inference tractable.

\end{frame}

\begin{frame}{Maximum-likelihood estimator in PPCA}

At the ML optimum, one may write
\[
W_{\mathrm{ML}} = U_k (\Lambda_k - \sigma^2 I)^{1/2} R,
\]
where \(R\) is orthogonal.
Thus the column space of \(W_{\mathrm{ML}}\) is the PCA subspace.
So PPCA preserves the principal directions while adding an explicit isotropic noise model.

\end{frame}

\begin{frame}{Deriving the ELBO carefully}

Insert any density \(q_\phi(z\mid x)\):
\[
\log p_\theta(x)
=
\log \int q_\phi(z\mid x)\frac{p_\theta(x,z)}{q_\phi(z\mid x)}\,dz.
\]
Applying Jensen:
\[
\log p_\theta(x)
\ge
\int q_\phi(z\mid x)\log \frac{p_\theta(x,z)}{q_\phi(z\mid x)}\,dz.
\]
This is the fundamental inequality underlying VAEs.

\end{frame}

\begin{frame}{KL decomposition identity}

Starting from Bayes' rule,
\[
p_\theta(z\mid x)=\frac{p_\theta(x,z)}{p_\theta(x)},
\]
one shows
\[
\log p_\theta(x)
=
\mathcal{L}_{\mathrm{ELBO}}(x)
+
\mathrm{KL}\bigl(q_\phi(z\mid x)\|p_\theta(z\mid x)\bigr).
\]
Thus maximizing the ELBO simultaneously:
\begin{itemize}
\item increases the likelihood,
\item improves approximate inference.
\end{itemize}

\end{frame}

\begin{frame}{Gaussian VAE components}

A common VAE parameterization is
\[
q_\phi(z\mid x)=\mathcal{N}(\mu_\phi(x), \diag(\sigma_\phi(x)^2)),
\]
\[
p(z)=\mathcal{N}(0,I),
\]
\[
p_\theta(x\mid z)=\mathcal{N}(\mu_\theta(z), \sigma_x^2 I)
\]
or a Bernoulli model for binary data.
The analytic KL term is especially convenient for Gaussian encoder and prior.

\end{frame}

\begin{frame}{Closed form KL for Gaussian latent variables}

If
\[
q(z\mid x)=\mathcal{N}(\mu,\diag(\sigma^2)),
\qquad
p(z)=\mathcal{N}(0,I),
\]
then
\[
\mathrm{KL}(q\|p)
=
\frac{1}{2}\sum_{i=1}^k \left(\mu_i^2+\sigma_i^2-\log \sigma_i^2 -1\right).
\]
This makes VAE training computationally practical.

\end{frame}

\begin{frame}{Posterior collapse}

A known pathology in VAEs is posterior collapse:
\[
q_\phi(z\mid x) \approx p(z).
\]
Then the latent variable carries little information about \(x\).
This is more likely when:
\begin{itemize}
\item the decoder is too powerful,
\item KL regularization dominates,
\item optimization is poorly balanced.
\end{itemize}

\end{frame}

\begin{frame}{Information bottleneck formulation}

The information bottleneck objective is often written schematically as
\[
\min I(X;Z) - \beta I(Z;Y),
\]
or, in unsupervised settings, as a tradeoff between compression and faithful representation.
The VAE KL term controls \(I(X;Z)\) indirectly by encouraging the aggregate posterior to remain close to the prior.

\end{frame}

\begin{frame}{Rate--distortion and beta-VAE}

In \(\beta\)-VAE,
\[
\mathcal{L}_\beta
=
\mathbb{E}_{q_\phi(z\mid x)}[\log p_\theta(x\mid z)]
-\beta \,\mathrm{KL}(q_\phi(z\mid x)\|p(z)).
\]
Larger \(\beta\) increases compression but often sacrifices reconstruction quality.
This mirrors classical rate--distortion theory: lower rate comes at higher distortion.

\end{frame}

\begin{frame}{From PCA to VAE as a hierarchy}

The conceptual chain is:
\[
\text{PCA}
\quad\to\quad
\text{Linear autoencoder}
\quad\to\quad
\text{Probabilistic PCA}
\quad\to\quad
\text{Variational autoencoder}.
\]
Each step adds:
\begin{itemize}
\item more flexibility,
\item more probabilistic structure,
\item richer latent geometry.
\end{itemize}

\end{frame}

\begin{frame}{Lecture 4 exercise}

Show that if the decoder is linear and Gaussian and the posterior family is exact, then the latent-variable model reduces to a linear-Gaussian inference problem.
How does this explain the relation between VAEs and PPCA?

\end{frame}

\end{document}
